The eigensheaf problem originates in Drinfeld's rank-two construction and
Laumon's geometric formulation for general linear groups
\cite{Drinfeld,Laumon}. Frenkel, Gaitsgory, and Vilonen reduced unramified
$\GL_n$ eigensheaf existence to a vanishing conjecture \cite{FGV}, which
Gaitsgory subsequently proved \cite{GaitsgoryVanishing}. Geometric Satake
identifies the local Hecke kernels with representations of the dual group
and supplies their tensor and fusion structure \cite{MV}. These results
established the geometric form of the passage from a prescribed local
system to an automorphic sheaf.

Our immediate antecedent is \cite[Theorem~1.1]{CT}, whose main theorem
has one marked point and regular-unipotent monodromy for a broader class
of groups. Here the group is $\SL_n$ with $p>n$, there are several marked
points, and every nilpotent boundary logarithm is allowed. The
two-specialization construction and the local arguments follow that
companion; we prove the simultaneous level geometry and compatibility
arguments needed for this extension.

For semisimple groups over $\mathbb C$, Beilinson and Drinfeld constructed
eigensheaves from opers by quantizing the Hitchin system \cite[\S0.2]{BD}.
The recent five-paper collaboration proves categorical forms of the
unramified correspondence in characteristic zero \cite{GLCI,GLCV}.
Arinkin, Gaitsgory, Kazhdan, Raskin, Rozenblyum, and Varshavsky introduced
the stack of local systems with restricted variation and its action on
sheaves with nilpotent singular support \cite{AGKRRV}.
Gaitsgory and Raskin use this framework to formulate the geometric adic
correspondence and compare characteristic zero with positive characteristic
\cite{GR}. The input needed here is the characteristic-zero restricted
equivalence \cite[Main Theorem~1.3.9(ii)]{GR}, through the unramified
eigencomplex construction of \cite[Lemma~8.1]{CT}. Thus the categorical
theory enters on a smooth projective curve before level structures are
introduced.

Nilpotent singular support links this input to ramification. The geometry
of the global nilpotent cone was developed by Laumon, Faltings, and Ginzburg
\cite{LaumonNilpotent,Faltings,GinzburgNilpotent}. Beilinson's theory of
singular support for constructible \'etale sheaves, and Barrett's extension
to adic coefficients, give the function tests used below
\cite{Beilinson,Barrett}. The Hecke detection argument in
\cite[\S\S20.5--20.8, Appendix~H]{AGKRRV} uses a cocharacter in the center
of a Levi subgroup and an isolated cotangent intersection. The field and
specialization arguments of \cite[\S\S3--4]{CT} adapt that geometry. We establish
the versions required for all the level stacks occurring here.
Geometric nearby cycles and their perverse exactness rest on
Hansen--Scholze and Gaitsgory--Raskin \cite{HS,GR}.

The characteristic-zero level arguments use Betti sheaf theory.
Nadler and Yun prove local constancy of moving Hecke operators on the
nilpotent category and construct its spectral action, already allowing
finitely many level points \cite[Theorems~6.2.2 and~6.3.7]{NY}.
Their automorphic gluing construction relates twisted nodal curves to
smoothings \cite{NYGluing}; its description of bundle types and branch
gluing is the geometric antecedent of our degeneration. Our proof also
needs a separate nonvanishing statement for the particular nearby cycles
being taken. At a puncture, the local spectral input comes from
Gaitsgory's central sheaves \cite{GaitsgoryCentral}, the affine-flag
equivalence of Arkhipov and Bezrukavnikov \cite{AB}, and Dhillon and
Taylor's extension to the universal monodromic setting \cite{DT}.
The last supplies the local kernels used in the descent argument of
\cite[\S7]{CT}, which we globalize with all the other levels retained.
Dhillon and Taylor's sequel proves the tame local Betti correspondence
\cite{DTLocal}.

Several-point ramification has earlier geometric realizations. Uit de Bos
gives an explicit $\GL_2$ correspondence for the projective line with four
punctures and pure irreducible parameters with unipotent monodromy
\cite{UitDeBos}. Shen proves a Hecke-compatible equivalence over an open
spectral locus for parabolic $\GL_n$ in positive characteristic, using
crystalline $D$-modules \cite{Shen}. In the characteristic-zero de Rham
setting, F{\ae}rgeman gives a construction for irreducible regular-singular
parameters conditional on the forthcoming factorization input and local
compatibility conjecture specified in that work
\cite[Theorem~1.4.1.1, \S\S1.3.7, 1.4.3, Conjecture~4.2.1]{Faergeman}.
In the cited version,
regular holonomicity and generic perversity in the ramified setting are
stated as expectations \cite[Remark~1.4.1.2]{Faergeman}. The present
theorem concerns geometric adic sheaves in positive characteristic and
includes local constructibility and perversity.
